\section{Introduction}
\label{sec:intro}

The modern paradigm of deep learning has undergone a monumental shift toward modularity. Rather than training colossal, monolithic neural networks for every individual application, the standard practice fine-tunes specialized expert pathways via Parameter-Efficient Fine-Tuning (PEFT) techniques, such as Low-Rank Adaptation (LoRA) \cite{hu21lora}, on top of frozen pre-trained backbones \cite{dosovitskiy20vit, devlin18bert}. This modularity is particularly vital for resource-constrained edge devices (e.g., mobile phones, autonomous robots, and smart sensors), which must rapidly adapt to multi-task workloads. In real-world deployment, these systems are subjected to highly heterogeneous and noisy input streams, where the underlying task identity shifts sample-by-sample (e.g., transitioning from digit classification to object detection to semantic segmentation in a fraction of a second).

To address this challenge without the prohibitive storage and memory footprint of hosting multiple standalone models, researchers have turned to \emph{model merging} \cite{wortsman22soups}. Static weight-space model merging methods, such as Task Arithmetic \cite{ilharco22taskarithmetic}, TIES-Merging \cite{yadav23ties}, and DARE \cite{yu23dare}, average task-specific parameter vectors into a single monolithic weight tensor. While computationally attractive at inference time, static merging is fundamentally a lossy process. When faced with highly diverse, mixed-task streaming batches, static merged models catastrophically fail due to \emph{Heterogeneity Collapse}, where conflicting expert parameters neutralize each other, degrading performance across all tasks.

To prevent this collapse, dynamic activation-space ensembling methods, such as SABLE \cite{liang23sable} and SPS-ZCA \cite{zhang24spszca}, blend expert representations dynamically sample-wise inside the forward pass using sample-dependent routing coefficients. Despite their improvements, these methods rely on a foundational assumption that we challenge in this work: \textbf{they treat the sequence of layers $l = 1 \dots L$ as decoupled, independent execution blocks}. In SABLE and SPS-ZCA, routing coefficients are either computed once globally at an early layer or computed independently per layer without memory of the preceding routing states. 

From a dynamical systems perspective, this stateless execution lacks physical continuity. When a model processes highly heterogeneous stream samples under noisy conditions, stateless routers exhibit sharp, discontinuous switching spikes and high-frequency "routing weight jitter" between adjacent layers. Biological brains and complex physical systems never transition states instantaneously; they obey continuous, smooth kinetics governed by physical memory (inertia) and thermodynamic constraints. Furthermore, systems-level scheduling wrappers like Micro-Batch Homogenization (MBH) attempt to restore representational stability by grouping samples, but they require $O(K)$ sequential backbone passes. This introduces a severe $4\times$ latency penalty on edge hardware, violating the core requirement of real-time streaming adaptation.

\begin{figure*}[t]
  \centering
  \begin{subfigure}[b]{0.32\textwidth}
    \centering
    \includegraphics[width=\textwidth]{results/fig1.png}
    \caption{Performance Sweep across seeds}
    \label{fig:sweep}
  \end{subfigure}
  \hfill
  \begin{subfigure}[b]{0.32\textwidth}
    \centering
    \includegraphics[width=\textwidth]{results/batch_size_heterogeneity.png}
    \caption{Heterogeneity Robustness}
    \label{fig:heterogeneity}
  \end{subfigure}
  \hfill
  \begin{subfigure}[b]{0.32\textwidth}
    \centering
    \includegraphics[width=\textwidth]{results/layer_trajectory.png}
    \caption{Concentration Trajectories}
    \label{fig:trajectories}
  \end{subfigure}
  \caption{\textbf{[SIMULATED/SANDBOX]} Overview of ChemMerge performance and trajectories. (a) Performance sweep across 10 random seeds shows ChemMerge closely tracking the Expert Ceiling and outperforming nearest-centroid baselines while remaining highly competitive with other ensembling baselines. (b) Accuracy under varying batch heterogeneity ($B=1$ to $256$) demonstrates ChemMerge's complete immunity to collapse. (c) Layer-wise chemical concentration trajectories show smooth, continuous state transitions across deep layers, physically filtering high-frequency routing jitter.}
  \label{fig:main_plots}
\end{figure*}

In this work, we present a radical, high-risk conceptual pivot that bridges systems biochemistry and deep neural networks. We propose \textbf{ChemMerge}, a training-free, continuous-time paradigm that models the flow of representations through the depth of a deep neural network as a multi-component chemical reactor governed by non-equilibrium reaction kinetics. By viewing task-specific LoRA experts as reactive chemical species, pre-computed task centroids as catalytic enzymes, and hidden activations as a reacting solution, we introduce physical temporal memory across sequential layers.

Rather than evaluating stateless similarities, ChemMerge tracks a continuous, sample-wise expert concentration state vector $C_b^{(l)} \in [0, 1]^K$ that evolves across the depth of the network. The state transitions are governed by first-order reversible chemical kinetics, where the forward reaction rate is determined by the catalytic affinity of the input to the expert's centroid (via a temperature-scaled Arrhenius equation), and the backward decay rate maintains representation plasticity. Discretizing this system via an explicit Euler step across sequential layers physically dampens high-frequency routing jitter and prevents sequential representation drift, delivering a stable and smooth representation trajectory.

While our continuous-time formulation is grounded in the elegant mathematics of systems biochemistry, we highlight a profound mathematical duality: the discretized kinetics governing ChemMerge are mathematically equivalent to a \emph{state-dependent, adaptive Exponential Moving Average (EMA) low-pass filter} (or first-order digital smoothing filter) applied directly to the routing coefficients across depth. Rather than applying a heuristic, static filter that introduces representational lag, ChemMerge's chemical framework provides a principled, state-dependent mechanism where the local smoothing rate adapts dynamically to the input similarity. This synthesis of continuous-depth physical kinetics with adaptive digital signal processing (DSP) provides a robust, mathematically rigorous foundation that resolves the long-standing stability-accuracy trade-off in streaming model serving.

Our core contributions are:
\begin{itemize}
    \item \textbf{A Novel Biochemical Paradigm for Deep Architectures:} We completely reject the standard stateless layer-wise decoupled ensembling framework. We establish a mathematically sound, training-free formulation of deep neural inference as a non-equilibrium chemical reaction network.
    \item \textbf{Continuous State Tracking via Kinetic Routing:} We introduce Non-Equilibrium Kinetic Routing (NEKR), which maintains physical continuity and inertia across sequential layers via first-order kinetics differential equations, resolving the layer-to-layer routing jitter paradox.
    \item \textbf{Evaluation in a High-Fidelity Simulation Sandbox:} Evaluated inside a 14-layer, 192-dimensional Analytical Coordinate Sandbox (ICS) across 10 independent random seeds (simulating MNIST, Fashion-MNIST, CIFAR-10, and SVHN streams), ChemMerge achieves a Joint Mean of \textbf{78.11\%} (homogeneous) and \textbf{78.06\%} (heterogeneous), recovering \textbf{98.81\%} of the Expert Ceiling. This sandbox isolates ensembling and routing mechanics with high scientific transparency.
    \item \textbf{Routing-Only Validation on Foundation Models:} On a pre-trained Vision Transformer ($\text{ViT-B/16}$), we evaluate the routing trajectories on real-world activation manifolds. ChemMerge reduces layer-to-layer ensembling weight routing jitter by up to \textbf{9.9$\times$} compared to stateless routing and over \textbf{2.15$\times$} compared to SABLE (under identical sensitivities). While absolute classification accuracies are statistically comparable due to overlapping standard deviations, ChemMerge successfully resolves the accuracy-stability trade-off with constant $O(1)$ edge serving latency.
    \item \textbf{Robustness to Structural Collapses:} We empirically demonstrate that ChemMerge is completely immune to both Heterogeneity Collapse ($B=256$) and Vectorization Collapse ($B=1$), resolving the critical streaming adaptivity bottleneck.
\end{itemize}
